\chapter{Auth Middleware for Protected APIs}

Chapters 18 and 19 covered issuing and storing a JWT. This chapter puts it
to work: a reusable Express middleware that stands in front of any route
that should only be reachable by a logged-in user, verifies the token, and
loads the actual user document so controllers can trust \texttt{req.user}.

\begin{diagrambox}[Protected Request Flow]
\begin{verbatim}
 Request (with "token" cookie)
        |
        v
 Auth Middleware (protect)
        |
        v
 JWT Verification (jwt.verify against JWT_SECRET)
        |
        v
 User Lookup (User.findById(decoded.id))
        |
        v
 req.user = user document (password excluded)
        |
        v
 next() -> Controller (trusts req.user, no re-checking)
\end{verbatim}
\end{diagrambox}

\section{The \texttt{protect} Middleware}

\begin{lstlisting}[language=JavaScript]
// middleware/authMiddleware.js
import jwt from "jsonwebtoken";
import User from "../models/User.js";

export const protect = async (req, res, next) => {
  const token = req.cookies.token;

  if (!token) {
    return res.status(401).json({ success: false, message: "Not authenticated" });
  }

  try {
    const decoded = jwt.verify(token, process.env.JWT_SECRET);

    const user = await User.findById(decoded.id).select("-password");
    if (!user) {
      return res.status(401).json({ success: false, message: "User no longer exists" });
    }

    req.user = user;
    next();
  } catch (err) {
    return res.status(401).json({ success: false, message: "Invalid or expired token" });
  }
};
\end{lstlisting}

\subsection{Walking Through Each Failure Case}

\begin{center}
\begin{tabularx}{\textwidth}{lX}
\toprule
\textbf{Scenario} & \textbf{Result} \\
\midrule
No cookie sent at all & \texttt{req.cookies.token} is \texttt{undefined} -- immediate 401, never touches \texttt{jwt.verify}. \\
Token tampered with or signed by a different secret & \texttt{jwt.verify} throws a \texttt{JsonWebTokenError} -- caught, 401. \\
Token past its \texttt{exp} claim & \texttt{jwt.verify} throws a \texttt{TokenExpiredError} -- caught, 401. \\
Token valid, but user was deleted after it was issued & \texttt{jwt.verify} succeeds, \texttt{User.findById} returns \texttt{null} -- explicit 401. \\
Token valid, user exists & \texttt{req.user} is set, \texttt{next()} passes control to the controller. \\
\bottomrule
\end{tabularx}
\end{center}

\begin{notebox}[NOTE]
\texttt{.select("-password")} excludes the hashed password from the document
attached to \texttt{req.user}. Even though it's a hash and not the plain
password, there's no reason for it to ever leave the database layer.
\end{notebox}

\section{Using the Middleware}

Apply \texttt{protect} to any route that requires a logged-in user, simply
by inserting it before the controller in the route definition:

\begin{lstlisting}[language=JavaScript]
// routes/taskRoutes.js
import express from "express";
import { protect } from "../middleware/authMiddleware.js";
import { getTasks, createTask } from "../controllers/taskController.js";

const router = express.Router();

router.get("/tasks", protect, getTasks);
router.post("/tasks", protect, createTask);

export default router;
\end{lstlisting}

Inside \texttt{getTasks}, the controller can now safely reference
\texttt{req.user} without re-verifying anything:

\begin{lstlisting}[language=JavaScript]
export const getTasks = async (req, res, next) => {
  try {
    const tasks = await Task.find({ user: req.user._id });
    res.status(200).json({ success: true, data: tasks });
  } catch (err) {
    next(err);
  }
};
\end{lstlisting}

\begin{tipbox}[TIP]
Express runs middleware in the order listed. Because \texttt{protect}
appears before \texttt{getTasks} in the route definition, it always executes
first -- if it doesn't call \texttt{next()}, the controller never runs. This
is what makes the route "protected."
\end{tipbox}

\whatwhyhow
{An Express middleware that reads the JWT cookie, verifies it, loads the corresponding user, and attaches it to \texttt{req.user} before letting the request continue.}
{Controllers need a trustworthy, verified identity without each one re-implementing token verification and database lookups themselves.}
{Insert \texttt{protect} as a route parameter before the controller function: \texttt{router.get("/tasks", protect, getTasks)}; any missing, invalid, or expired token short-circuits with a 401 before the controller ever runs.}
